% Generated draft: requires mathematical and source review.
\documentclass[11pt]{article}
\usepackage[margin=25mm]{geometry}
\usepackage{fontspec}
\setmainfont{Latin Modern Roman}
\usepackage{amsmath,amssymb,mathtools,unicode-math}
\setmathfont{Latin Modern Math}
\usepackage{enumitem,needspace,xurl,hyperref}
\setlength{\parindent}{0pt}
\setlength{\parskip}{5pt}
\setlength{\emergencystretch}{3em}
\newcommand{\N}{\mathbb{N}}
\newcommand{\Z}{\mathbb{Z}}
\newcommand{\Q}{\mathbb{Q}}
\newcommand{\R}{\mathbb{R}}
\newcommand{\C}{\mathbb{C}}
\newcommand{\F}{\mathbb{F}}
\newcommand{\A}{\mathbb{A}}
\newcommand{\PP}{\mathbb{P}}
\newcommand{\E}{\mathbb{E}}
\newcommand{\eps}{\varepsilon}
\newcommand{\dd}{\,\mathrm{d}}
\begin{document}
\section*{UnsolvedMath ID: 30001142}
\textbf{Source problem code:} catalogue entry
\subsection{Definitions and mathematical statement}
Let \(S^2=\{x\in\R^3:\|x\|=1\}\), with its usual smooth structure. A \(C^\infty\) embedding \(f:S^2\to\R^3\) is a smooth injective immersion that is a homeomorphism onto its image. Six points form the vertices of a nondegenerate regular octahedron if they are precisely
\[
\{c+r u_1,c-r u_1,c+r u_2,c-r u_2,c+r u_3,c-r u_3\},
\]
where \(c\in\R^3\), \(r>0\), and \((u_1,u_2,u_3)\) is an orthonormal basis of \(\R^3\).

\textbf{Question.} For every \(C^\infty\) embedding \(f:S^2\to\R^3\), do there exist \(c\in\R^3\), \(r>0\), and an orthonormal basis \((u_1,u_2,u_3)\) such that
\[
 c+r u_i\in f(S^2)\quad\text{and}\quad c-r u_i\in f(S^2)
 \qquad\text{for every }i\in\{1,2,3\}?
\]
Only the six vertices are required to lie on the embedded sphere; containment of the edges, faces, or solid octahedron in the surface is not required. The supplied background also mentions metric two-spheres, but gives no definition or hypotheses for that broader formulation; the question here retains the explicitly stated smooth embedded setting.
\subsection{Short English statement}
Does every smoothly embedded two-sphere in \(\R^3\) contain the six vertices of a nondegenerate regular octahedron? The supplied record labels this problem open; that status is supplied curation, not an independently verified conclusion.
\subsection{Sources}
\begin{itemize}
\item Local UnsolvedMath registry, permanent ID \(30001142\), title: \emph{Regular Octahedra Inscribed in Smooth Two-Spheres}; source problem code \texttt{catalogue entry}. The record attributes the problem to the Oberwolfach Reports collection and gives proposed year 2009, without naming an individual proposer.
\item \emph{Discrete Differential Geometry} (2009), Oberwolfach Reports, pp.\ 75--144, DOI 10.4171/owr/2009/02. Supplied source report and literature-source link: \url{https://doi.org/10.4171/owr/2009/02}. The supplied background describes these reports as documenting mathematical workshops and their open-problem discussions.
\item Matschke, \emph{square peg problem and relatives}, as labeled in the supplied literature sources: \url{https://doi.org/10.1090/tran/8649}. The supplied assessment says that the regular-octahedron problem for smooth or metric two-spheres remains open and that Matschke identifies the dimension-three crosspolytope case as open after finding a topological obstruction. This assessment is supplied dated curation checked on \textbf{2026-08-21}, not independently verified here; its status and cited claims require independent verification.
\item The registry marks the statement \texttt{corrected\_verified} and describes its processing as meaning-preserving semantic, typographic, and LaTeX cleanup without unsupported source-dependent additions. These are supplied provenance labels, not an independent mathematical or source verification of this draft.
\end{itemize}
\end{document}
